\documentclass[aspectratio=169,10pt]{beamer}
% ---------------------------------------------------------------
% R4.02 Qualite de developpement -- preambule commun des diapositives
% Charte reprise de R3.07 (metropolis, dark tech)
% Palette IUT partagee (violet du logo + accent or)
% Utilisation :
%   \documentclass[aspectratio=169,10pt]{beamer}
%   % ---------------------------------------------------------------
% R4.02 Qualite de developpement -- preambule commun des diapositives
% Charte reprise de R3.07 (metropolis, dark tech)
% Palette IUT partagee (violet du logo + accent or)
% Utilisation :
%   \documentclass[aspectratio=169,10pt]{beamer}
%   % ---------------------------------------------------------------
% R4.02 Qualite de developpement -- preambule commun des diapositives
% Charte reprise de R3.07 (metropolis, dark tech)
% Palette IUT partagee (violet du logo + accent or)
% Utilisation :
%   \documentclass[aspectratio=169,10pt]{beamer}
%   \input{../../latex/common/r402-slides-preamble.tex}
%   \graphicspath{{../../latex/common/assets/}{img/}{figs/}}
% ---------------------------------------------------------------
\usetheme{metropolis}
\usepackage{appendixnumberbeamer}

\usepackage[T1]{fontenc}
\usepackage[utf8]{inputenc}
\usepackage[french]{babel}
\usepackage[scale=0.94]{tgheros}
\renewcommand{\familydefault}{\sfdefault}
\usepackage{amsmath,amssymb,amsfonts}
\usepackage{graphicx}
\usepackage{booktabs}
\usepackage{multirow}
\usepackage{tikz}
\usetikzlibrary{shapes.geometric, shapes, arrows, arrows.meta, positioning, calc, fit}
\usepackage{xcolor}
\usepackage{listings}
\usepackage[skins,listings]{tcolorbox}
\usepackage{qrcode}
\usepackage{eurosym}
\usepackage{hyperref}

% ---------------------------------------------------------------
% PALETTE (dark-mode tech: near-black canvas, violet (logo IUT) + gold accents)
% ---------------------------------------------------------------
\definecolor{qabg}{HTML}{0E1016}      % canvas background
\definecolor{qapanel}{HTML}{181B24}   % card / block panel
\definecolor{qapanelalt}{HTML}{20242F}% secondary panel tint
\definecolor{qaline}{HTML}{2B303D}    % hairline rules / table rules
\definecolor{qatext}{HTML}{ECEEF2}    % primary text (off-white)
\definecolor{qamuted}{HTML}{8B93A6}   % secondary / caption text
\definecolor{qaaccent}{HTML}{D93ED7}  % primary accent -- violet du logo IUT, eclairci pour fond sombre
\definecolor{qaaccent2}{HTML}{F2B84B} % secondary accent -- or/ambre
\definecolor{qacodekw}{HTML}{C792EA}  % code : mots-cles
\definecolor{qacodestr}{HTML}{C3E88D} % code : chaines
\definecolor{qacodecom}{HTML}{7F8AA3} % code : commentaires
\definecolor{qacodeann}{HTML}{F2B84B} % code : annotations / javadoc

\setbeamercolor{background canvas}{bg=qabg}
\setbeamercolor{normal text}{fg=qatext}
\setbeamercolor{palette primary}{bg=qabg, fg=qamuted}
\setbeamercolor{frametitle}{bg=qabg, fg=qatext}
\setbeamercolor{progress bar}{fg=qaaccent, bg=qaline}
\setbeamercolor{alerted text}{fg=qaaccent2}
\setbeamercolor{block title}{bg=qapanel, fg=qaaccent}
\setbeamercolor{block body}{bg=qapanel, fg=qatext}
\setbeamercolor{block title alerted}{bg=qapanel, fg=qaaccent2}
\setbeamercolor{block body alerted}{bg=qapanel, fg=qatext}
\setbeamercolor{block title example}{bg=qapanelalt, fg=qaaccent}
\setbeamercolor{block body example}{bg=qapanelalt, fg=qatext}
\setbeamercolor{example text}{fg=qaaccent}
\setbeamercolor{itemize item}{fg=qaaccent}
\setbeamercolor{itemize subitem}{fg=qaaccent2}
\setbeamercolor{enumerate item}{fg=qaaccent}
\setbeamercolor{enumerate subitem}{fg=qaaccent2}
\setbeamercolor{caption}{fg=qamuted}
\setbeamercolor{page number in head/foot}{fg=qamuted, bg=qabg}
\setbeamercolor{title}{fg=qatext}
\setbeamercolor{subtitle}{fg=qamuted}
\setbeamercolor{author}{fg=qatext}
\setbeamercolor{institute}{fg=qamuted}
\setbeamercolor{date}{fg=qamuted}
\setbeamercolor{section in toc}{fg=qaaccent}
\setbeamercolor{subsection in toc}{fg=qamuted}
% Sommaire de debut de section : section en cours en violet et gras, les autres en gris
\setbeamercolor{section in toc shaded}{fg=qamuted}
\setbeamercolor{subsection in toc shaded}{fg=qamuted}

\setbeamertemplate{blocks}[rounded][shadow=false]
\setbeamertemplate{itemize item}{\raisebox{0.15em}{\scriptsize\textcolor{qaaccent}{\ensuremath{\blacktriangleright}}}}
\setbeamertemplate{itemize subitem}{\raisebox{0.1em}{\tiny\textcolor{qaaccent2}{\ensuremath{\bullet}}}}

\setbeamerfont{frametitle}{series=\bfseries, size=\Large}
\setbeamerfont{title}{series=\bfseries, size=\Large}
\setbeamerfont{subtitle}{size=\normalsize}
\setbeamerfont{institute}{size=\scriptsize}
\setbeamerfont{date}{size=\scriptsize}
\setbeamerfont{author}{size=\small}

\metroset{sectionpage=none, subsectionpage=none, numbering=fraction, progressbar=frametitle}
\setbeamertemplate{section in toc}[sections numbered]

% ---------------------------------------------------------------
% LISTINGS -- theme sombre (lstlisting reste utilisable tel quel)
% \lstcode : police + couleur de base, a reutiliser dans les
%            surcharges locales (ex. basicstyle=\lstcode\scriptsize)
% ---------------------------------------------------------------
\newcommand{\lstcode}{\ttfamily\color{qatext}}

\lstdefinelanguage{TypeScript}{
  keywords={typeof, new, true, false, catch, function, return, null, switch, var,
    if, in, while, do, else, case, break, const, let, interface, class, extends,
    implements, enum, export, import, this, super, constructor, public, private,
    protected, static, async, await, throw, throws, try, finally, void},
  ndkeywords={boolean, any, number, string, int, long, double},
  sensitive=true,
  comment=[l]{//},
  morecomment=[s]{/*}{*/},
  morestring=[b]',
  morestring=[b]"
}

\lstset{
  language=Java,
  basicstyle=\lstcode\footnotesize,
  keywordstyle=\color{qacodekw}\bfseries,
  ndkeywordstyle=\color{qaaccent2}\bfseries,
  identifierstyle=\color{qatext},
  stringstyle=\color{qacodestr},
  commentstyle=\color{qacodecom}\itshape,
  morecomment=[s][\color{qacodeann}]{/**}{*/},
  backgroundcolor=\color{qapanelalt},
  rulecolor=\color{qaline},
  frame=single,
  framesep=3pt,
  xleftmargin=4pt,
  xrightmargin=4pt,
  keepspaces=true,
  columns=fullflexible,
  showstringspaces=false,
  breaklines=true,
  tabsize=2,
  escapeinside={(*@}{@*)},
  % accents UTF-8 dans le code (listings ne gere pas l'UTF-8 multi-octets)
  literate={é}{{\'e}}1 {è}{{\`e}}1 {ê}{{\^e}}1 {à}{{\`a}}1 {â}{{\^a}}1
           {ù}{{\`u}}1 {û}{{\^u}}1 {ô}{{\^o}}1 {î}{{\^i}}1 {ç}{{\c c}}1
           {É}{{\'E}}1 {À}{{\`A}}1 {€}{{\euro{}}}1 {…}{{\ldots}}1,
}

% ---------------------------------------------------------------
% CARD-STYLE BOXES
% ---------------------------------------------------------------
\tcbset{
  figcard/.style={colback=white, colframe=white, boxrule=0pt, arc=4pt,
    left=3pt, right=3pt, top=3pt, bottom=3pt, boxsep=0pt,
    halign=center, valign=center},
  datacard/.style={colback=qapanel, colframe=qaline, boxrule=0.4pt, arc=3pt,
    left=8pt, right=8pt, top=6pt, bottom=6pt, coltext=qatext,
    before upper={\color{qatext}}, halign=center},
  codebox/.style={colback=qapanelalt, colframe=qaline, boxrule=0.4pt, arc=3pt,
    left=8pt, right=8pt, top=5pt, bottom=5pt, listing only,
    listing options={basicstyle=\ttfamily\small\color{qaaccent},
      breaklines=true, showstringspaces=false, columns=fullflexible,
      keepspaces=true, frame=none, backgroundcolor={}}}
}
\newtcolorbox{figcard}[1][]{figcard,#1}
\newtcolorbox{datacard}{datacard}
\newtcblisting{codebox}{codebox}

% ---------------------------------------------------------------
% HANDY MACROS
% ---------------------------------------------------------------
% Mention d'elaboration (discrete : plus sombre que le texte secondaire)
\definecolor{qaaimention}{HTML}{4E5566}
\newcommand{\raimention}{{\tiny\color{qaaimention}Supports \'elabor\'es avec l'aide de Claude Code Opus 5.5}}
\newcommand{\hl}[1]{\alert{#1}}
\newcommand{\source}[1]{{\scriptsize\color{qamuted}#1}}

% Sommaire automatique en debut de section
\AtBeginSection[]{%
  \begin{frame}<beamer>{Plan du cours}
    \tableofcontents[currentsection]
  \end{frame}}

% Page de titre : \rtitleslide{numero du cours}{titre du cours}
% Les auteurs sont repris de \rauthors (modifiable avant \begin{document})
\newcommand{\rauthors}{Chouki Tibermacine}
\newcommand{\rtitleslide}[2]{%
\begin{frame}[plain,noframenumbering]
  \begin{center}
    \vspace{2mm}
    {\color{qaaccent}\rule{2.1cm}{1.2pt}}\\[6mm]
    {\Large\bfseries\color{qatext} R4.02}\\[1mm]
    {\Large\bfseries\color{qatext} Qualit\'e de d\'eveloppement}\\[5mm]
    {\large\color{qamuted} Cours #1}\\[1mm]
    {\large\color{qaaccent} #2}\\[6mm]
    \includegraphics[height=2.1cm]{iut_logo_transparent.png}

    \vspace{4mm}

    {\small\color{qatext} \rauthors}\\[1mm]
    {\scriptsize\color{qamuted} BUT 2 Informatique --- IUT de Vannes --- Universit\'e Bretagne Sud}\\[2mm]
    \raimention
  \end{center}
\end{frame}
}

% ---- Diapo de demonstration : \rdemo{numero}{titre}{contenu (items)}{fichier(s) a ouvrir} ----
\newcommand{\rdemo}[4]{%
\begin{frame}{D\'emo #1 --- #2}
  \begin{tikzpicture}[remember picture, overlay]
    \node[anchor=north east, fill=qaaccent2, text=qabg, font=\scriptsize\bfseries,
      rounded corners=3pt, inner sep=4pt] at ([xshift=-0.4cm, yshift=-0.35cm]current page.north east) {D\'EMO};
  \end{tikzpicture}
  \begin{columns}[T]
    \begin{column}{0.62\textwidth}
      \begin{itemize}
        #3
      \end{itemize}
    \end{column}
    \begin{column}{0.34\textwidth}
      \begin{datacard}
        \raggedright\scriptsize
        {\color{qamuted} Dans l'IDE}\\[1mm]
        #4
      \end{datacard}
    \end{column}
  \end{columns}
\end{frame}}

% ---- Diapo de fin « Questions ? » ----
\newcommand{\rcontact}{chouki.tibermacine@univ-ubs.fr}
\newcommand{\rquestionslide}{%
\begin{frame}{Questions ?}
  \begin{center}
    \vspace{2mm}
    \begin{tikzpicture}
      \node[regular polygon, regular polygon sides=6, draw=qaline, thick,
        minimum size=4.6cm, rotate=90] (hexin) {};
      \node[regular polygon, regular polygon sides=6, draw=qaaccent, dashed,
        thick, minimum size=6.2cm, rotate=90] (hexout) {};
      \node[font=\fontsize{92}{92}\selectfont\bfseries, color=qaaccent2] at (hexin) {?};
    \end{tikzpicture}\\[5mm]
    {\small\color{qamuted} \rcontact}
  \end{center}
\end{frame}}

%   \graphicspath{{../../latex/common/assets/}{img/}{figs/}}
% ---------------------------------------------------------------
\usetheme{metropolis}
\usepackage{appendixnumberbeamer}

\usepackage[T1]{fontenc}
\usepackage[utf8]{inputenc}
\usepackage[french]{babel}
\usepackage[scale=0.94]{tgheros}
\renewcommand{\familydefault}{\sfdefault}
\usepackage{amsmath,amssymb,amsfonts}
\usepackage{graphicx}
\usepackage{booktabs}
\usepackage{multirow}
\usepackage{tikz}
\usetikzlibrary{shapes.geometric, shapes, arrows, arrows.meta, positioning, calc, fit}
\usepackage{xcolor}
\usepackage{listings}
\usepackage[skins,listings]{tcolorbox}
\usepackage{qrcode}
\usepackage{eurosym}
\usepackage{hyperref}

% ---------------------------------------------------------------
% PALETTE (dark-mode tech: near-black canvas, violet (logo IUT) + gold accents)
% ---------------------------------------------------------------
\definecolor{qabg}{HTML}{0E1016}      % canvas background
\definecolor{qapanel}{HTML}{181B24}   % card / block panel
\definecolor{qapanelalt}{HTML}{20242F}% secondary panel tint
\definecolor{qaline}{HTML}{2B303D}    % hairline rules / table rules
\definecolor{qatext}{HTML}{ECEEF2}    % primary text (off-white)
\definecolor{qamuted}{HTML}{8B93A6}   % secondary / caption text
\definecolor{qaaccent}{HTML}{D93ED7}  % primary accent -- violet du logo IUT, eclairci pour fond sombre
\definecolor{qaaccent2}{HTML}{F2B84B} % secondary accent -- or/ambre
\definecolor{qacodekw}{HTML}{C792EA}  % code : mots-cles
\definecolor{qacodestr}{HTML}{C3E88D} % code : chaines
\definecolor{qacodecom}{HTML}{7F8AA3} % code : commentaires
\definecolor{qacodeann}{HTML}{F2B84B} % code : annotations / javadoc

\setbeamercolor{background canvas}{bg=qabg}
\setbeamercolor{normal text}{fg=qatext}
\setbeamercolor{palette primary}{bg=qabg, fg=qamuted}
\setbeamercolor{frametitle}{bg=qabg, fg=qatext}
\setbeamercolor{progress bar}{fg=qaaccent, bg=qaline}
\setbeamercolor{alerted text}{fg=qaaccent2}
\setbeamercolor{block title}{bg=qapanel, fg=qaaccent}
\setbeamercolor{block body}{bg=qapanel, fg=qatext}
\setbeamercolor{block title alerted}{bg=qapanel, fg=qaaccent2}
\setbeamercolor{block body alerted}{bg=qapanel, fg=qatext}
\setbeamercolor{block title example}{bg=qapanelalt, fg=qaaccent}
\setbeamercolor{block body example}{bg=qapanelalt, fg=qatext}
\setbeamercolor{example text}{fg=qaaccent}
\setbeamercolor{itemize item}{fg=qaaccent}
\setbeamercolor{itemize subitem}{fg=qaaccent2}
\setbeamercolor{enumerate item}{fg=qaaccent}
\setbeamercolor{enumerate subitem}{fg=qaaccent2}
\setbeamercolor{caption}{fg=qamuted}
\setbeamercolor{page number in head/foot}{fg=qamuted, bg=qabg}
\setbeamercolor{title}{fg=qatext}
\setbeamercolor{subtitle}{fg=qamuted}
\setbeamercolor{author}{fg=qatext}
\setbeamercolor{institute}{fg=qamuted}
\setbeamercolor{date}{fg=qamuted}
\setbeamercolor{section in toc}{fg=qaaccent}
\setbeamercolor{subsection in toc}{fg=qamuted}
% Sommaire de debut de section : section en cours en violet et gras, les autres en gris
\setbeamercolor{section in toc shaded}{fg=qamuted}
\setbeamercolor{subsection in toc shaded}{fg=qamuted}

\setbeamertemplate{blocks}[rounded][shadow=false]
\setbeamertemplate{itemize item}{\raisebox{0.15em}{\scriptsize\textcolor{qaaccent}{\ensuremath{\blacktriangleright}}}}
\setbeamertemplate{itemize subitem}{\raisebox{0.1em}{\tiny\textcolor{qaaccent2}{\ensuremath{\bullet}}}}

\setbeamerfont{frametitle}{series=\bfseries, size=\Large}
\setbeamerfont{title}{series=\bfseries, size=\Large}
\setbeamerfont{subtitle}{size=\normalsize}
\setbeamerfont{institute}{size=\scriptsize}
\setbeamerfont{date}{size=\scriptsize}
\setbeamerfont{author}{size=\small}

\metroset{sectionpage=none, subsectionpage=none, numbering=fraction, progressbar=frametitle}
\setbeamertemplate{section in toc}[sections numbered]

% ---------------------------------------------------------------
% LISTINGS -- theme sombre (lstlisting reste utilisable tel quel)
% \lstcode : police + couleur de base, a reutiliser dans les
%            surcharges locales (ex. basicstyle=\lstcode\scriptsize)
% ---------------------------------------------------------------
\newcommand{\lstcode}{\ttfamily\color{qatext}}

\lstdefinelanguage{TypeScript}{
  keywords={typeof, new, true, false, catch, function, return, null, switch, var,
    if, in, while, do, else, case, break, const, let, interface, class, extends,
    implements, enum, export, import, this, super, constructor, public, private,
    protected, static, async, await, throw, throws, try, finally, void},
  ndkeywords={boolean, any, number, string, int, long, double},
  sensitive=true,
  comment=[l]{//},
  morecomment=[s]{/*}{*/},
  morestring=[b]',
  morestring=[b]"
}

\lstset{
  language=Java,
  basicstyle=\lstcode\footnotesize,
  keywordstyle=\color{qacodekw}\bfseries,
  ndkeywordstyle=\color{qaaccent2}\bfseries,
  identifierstyle=\color{qatext},
  stringstyle=\color{qacodestr},
  commentstyle=\color{qacodecom}\itshape,
  morecomment=[s][\color{qacodeann}]{/**}{*/},
  backgroundcolor=\color{qapanelalt},
  rulecolor=\color{qaline},
  frame=single,
  framesep=3pt,
  xleftmargin=4pt,
  xrightmargin=4pt,
  keepspaces=true,
  columns=fullflexible,
  showstringspaces=false,
  breaklines=true,
  tabsize=2,
  escapeinside={(*@}{@*)},
  % accents UTF-8 dans le code (listings ne gere pas l'UTF-8 multi-octets)
  literate={é}{{\'e}}1 {è}{{\`e}}1 {ê}{{\^e}}1 {à}{{\`a}}1 {â}{{\^a}}1
           {ù}{{\`u}}1 {û}{{\^u}}1 {ô}{{\^o}}1 {î}{{\^i}}1 {ç}{{\c c}}1
           {É}{{\'E}}1 {À}{{\`A}}1 {€}{{\euro{}}}1 {…}{{\ldots}}1,
}

% ---------------------------------------------------------------
% CARD-STYLE BOXES
% ---------------------------------------------------------------
\tcbset{
  figcard/.style={colback=white, colframe=white, boxrule=0pt, arc=4pt,
    left=3pt, right=3pt, top=3pt, bottom=3pt, boxsep=0pt,
    halign=center, valign=center},
  datacard/.style={colback=qapanel, colframe=qaline, boxrule=0.4pt, arc=3pt,
    left=8pt, right=8pt, top=6pt, bottom=6pt, coltext=qatext,
    before upper={\color{qatext}}, halign=center},
  codebox/.style={colback=qapanelalt, colframe=qaline, boxrule=0.4pt, arc=3pt,
    left=8pt, right=8pt, top=5pt, bottom=5pt, listing only,
    listing options={basicstyle=\ttfamily\small\color{qaaccent},
      breaklines=true, showstringspaces=false, columns=fullflexible,
      keepspaces=true, frame=none, backgroundcolor={}}}
}
\newtcolorbox{figcard}[1][]{figcard,#1}
\newtcolorbox{datacard}{datacard}
\newtcblisting{codebox}{codebox}

% ---------------------------------------------------------------
% HANDY MACROS
% ---------------------------------------------------------------
% Mention d'elaboration (discrete : plus sombre que le texte secondaire)
\definecolor{qaaimention}{HTML}{4E5566}
\newcommand{\raimention}{{\tiny\color{qaaimention}Supports \'elabor\'es avec l'aide de Claude Code Opus 5.5}}
\newcommand{\hl}[1]{\alert{#1}}
\newcommand{\source}[1]{{\scriptsize\color{qamuted}#1}}

% Sommaire automatique en debut de section
\AtBeginSection[]{%
  \begin{frame}<beamer>{Plan du cours}
    \tableofcontents[currentsection]
  \end{frame}}

% Page de titre : \rtitleslide{numero du cours}{titre du cours}
% Les auteurs sont repris de \rauthors (modifiable avant \begin{document})
\newcommand{\rauthors}{Chouki Tibermacine}
\newcommand{\rtitleslide}[2]{%
\begin{frame}[plain,noframenumbering]
  \begin{center}
    \vspace{2mm}
    {\color{qaaccent}\rule{2.1cm}{1.2pt}}\\[6mm]
    {\Large\bfseries\color{qatext} R4.02}\\[1mm]
    {\Large\bfseries\color{qatext} Qualit\'e de d\'eveloppement}\\[5mm]
    {\large\color{qamuted} Cours #1}\\[1mm]
    {\large\color{qaaccent} #2}\\[6mm]
    \includegraphics[height=2.1cm]{iut_logo_transparent.png}

    \vspace{4mm}

    {\small\color{qatext} \rauthors}\\[1mm]
    {\scriptsize\color{qamuted} BUT 2 Informatique --- IUT de Vannes --- Universit\'e Bretagne Sud}\\[2mm]
    \raimention
  \end{center}
\end{frame}
}

% ---- Diapo de demonstration : \rdemo{numero}{titre}{contenu (items)}{fichier(s) a ouvrir} ----
\newcommand{\rdemo}[4]{%
\begin{frame}{D\'emo #1 --- #2}
  \begin{tikzpicture}[remember picture, overlay]
    \node[anchor=north east, fill=qaaccent2, text=qabg, font=\scriptsize\bfseries,
      rounded corners=3pt, inner sep=4pt] at ([xshift=-0.4cm, yshift=-0.35cm]current page.north east) {D\'EMO};
  \end{tikzpicture}
  \begin{columns}[T]
    \begin{column}{0.62\textwidth}
      \begin{itemize}
        #3
      \end{itemize}
    \end{column}
    \begin{column}{0.34\textwidth}
      \begin{datacard}
        \raggedright\scriptsize
        {\color{qamuted} Dans l'IDE}\\[1mm]
        #4
      \end{datacard}
    \end{column}
  \end{columns}
\end{frame}}

% ---- Diapo de fin « Questions ? » ----
\newcommand{\rcontact}{chouki.tibermacine@univ-ubs.fr}
\newcommand{\rquestionslide}{%
\begin{frame}{Questions ?}
  \begin{center}
    \vspace{2mm}
    \begin{tikzpicture}
      \node[regular polygon, regular polygon sides=6, draw=qaline, thick,
        minimum size=4.6cm, rotate=90] (hexin) {};
      \node[regular polygon, regular polygon sides=6, draw=qaaccent, dashed,
        thick, minimum size=6.2cm, rotate=90] (hexout) {};
      \node[font=\fontsize{92}{92}\selectfont\bfseries, color=qaaccent2] at (hexin) {?};
    \end{tikzpicture}\\[5mm]
    {\small\color{qamuted} \rcontact}
  \end{center}
\end{frame}}

%   \graphicspath{{../../latex/common/assets/}{img/}{figs/}}
% ---------------------------------------------------------------
\usetheme{metropolis}
\usepackage{appendixnumberbeamer}

\usepackage[T1]{fontenc}
\usepackage[utf8]{inputenc}
\usepackage[french]{babel}
\usepackage[scale=0.94]{tgheros}
\renewcommand{\familydefault}{\sfdefault}
\usepackage{amsmath,amssymb,amsfonts}
\usepackage{graphicx}
\usepackage{booktabs}
\usepackage{multirow}
\usepackage{tikz}
\usetikzlibrary{shapes.geometric, shapes, arrows, arrows.meta, positioning, calc, fit}
\usepackage{xcolor}
\usepackage{listings}
\usepackage[skins,listings]{tcolorbox}
\usepackage{qrcode}
\usepackage{eurosym}
\usepackage{hyperref}

% ---------------------------------------------------------------
% PALETTE (dark-mode tech: near-black canvas, violet (logo IUT) + gold accents)
% ---------------------------------------------------------------
\definecolor{qabg}{HTML}{0E1016}      % canvas background
\definecolor{qapanel}{HTML}{181B24}   % card / block panel
\definecolor{qapanelalt}{HTML}{20242F}% secondary panel tint
\definecolor{qaline}{HTML}{2B303D}    % hairline rules / table rules
\definecolor{qatext}{HTML}{ECEEF2}    % primary text (off-white)
\definecolor{qamuted}{HTML}{8B93A6}   % secondary / caption text
\definecolor{qaaccent}{HTML}{D93ED7}  % primary accent -- violet du logo IUT, eclairci pour fond sombre
\definecolor{qaaccent2}{HTML}{F2B84B} % secondary accent -- or/ambre
\definecolor{qacodekw}{HTML}{C792EA}  % code : mots-cles
\definecolor{qacodestr}{HTML}{C3E88D} % code : chaines
\definecolor{qacodecom}{HTML}{7F8AA3} % code : commentaires
\definecolor{qacodeann}{HTML}{F2B84B} % code : annotations / javadoc

\setbeamercolor{background canvas}{bg=qabg}
\setbeamercolor{normal text}{fg=qatext}
\setbeamercolor{palette primary}{bg=qabg, fg=qamuted}
\setbeamercolor{frametitle}{bg=qabg, fg=qatext}
\setbeamercolor{progress bar}{fg=qaaccent, bg=qaline}
\setbeamercolor{alerted text}{fg=qaaccent2}
\setbeamercolor{block title}{bg=qapanel, fg=qaaccent}
\setbeamercolor{block body}{bg=qapanel, fg=qatext}
\setbeamercolor{block title alerted}{bg=qapanel, fg=qaaccent2}
\setbeamercolor{block body alerted}{bg=qapanel, fg=qatext}
\setbeamercolor{block title example}{bg=qapanelalt, fg=qaaccent}
\setbeamercolor{block body example}{bg=qapanelalt, fg=qatext}
\setbeamercolor{example text}{fg=qaaccent}
\setbeamercolor{itemize item}{fg=qaaccent}
\setbeamercolor{itemize subitem}{fg=qaaccent2}
\setbeamercolor{enumerate item}{fg=qaaccent}
\setbeamercolor{enumerate subitem}{fg=qaaccent2}
\setbeamercolor{caption}{fg=qamuted}
\setbeamercolor{page number in head/foot}{fg=qamuted, bg=qabg}
\setbeamercolor{title}{fg=qatext}
\setbeamercolor{subtitle}{fg=qamuted}
\setbeamercolor{author}{fg=qatext}
\setbeamercolor{institute}{fg=qamuted}
\setbeamercolor{date}{fg=qamuted}
\setbeamercolor{section in toc}{fg=qaaccent}
\setbeamercolor{subsection in toc}{fg=qamuted}
% Sommaire de debut de section : section en cours en violet et gras, les autres en gris
\setbeamercolor{section in toc shaded}{fg=qamuted}
\setbeamercolor{subsection in toc shaded}{fg=qamuted}

\setbeamertemplate{blocks}[rounded][shadow=false]
\setbeamertemplate{itemize item}{\raisebox{0.15em}{\scriptsize\textcolor{qaaccent}{\ensuremath{\blacktriangleright}}}}
\setbeamertemplate{itemize subitem}{\raisebox{0.1em}{\tiny\textcolor{qaaccent2}{\ensuremath{\bullet}}}}

\setbeamerfont{frametitle}{series=\bfseries, size=\Large}
\setbeamerfont{title}{series=\bfseries, size=\Large}
\setbeamerfont{subtitle}{size=\normalsize}
\setbeamerfont{institute}{size=\scriptsize}
\setbeamerfont{date}{size=\scriptsize}
\setbeamerfont{author}{size=\small}

\metroset{sectionpage=none, subsectionpage=none, numbering=fraction, progressbar=frametitle}
\setbeamertemplate{section in toc}[sections numbered]

% ---------------------------------------------------------------
% LISTINGS -- theme sombre (lstlisting reste utilisable tel quel)
% \lstcode : police + couleur de base, a reutiliser dans les
%            surcharges locales (ex. basicstyle=\lstcode\scriptsize)
% ---------------------------------------------------------------
\newcommand{\lstcode}{\ttfamily\color{qatext}}

\lstdefinelanguage{TypeScript}{
  keywords={typeof, new, true, false, catch, function, return, null, switch, var,
    if, in, while, do, else, case, break, const, let, interface, class, extends,
    implements, enum, export, import, this, super, constructor, public, private,
    protected, static, async, await, throw, throws, try, finally, void},
  ndkeywords={boolean, any, number, string, int, long, double},
  sensitive=true,
  comment=[l]{//},
  morecomment=[s]{/*}{*/},
  morestring=[b]',
  morestring=[b]"
}

\lstset{
  language=Java,
  basicstyle=\lstcode\footnotesize,
  keywordstyle=\color{qacodekw}\bfseries,
  ndkeywordstyle=\color{qaaccent2}\bfseries,
  identifierstyle=\color{qatext},
  stringstyle=\color{qacodestr},
  commentstyle=\color{qacodecom}\itshape,
  morecomment=[s][\color{qacodeann}]{/**}{*/},
  backgroundcolor=\color{qapanelalt},
  rulecolor=\color{qaline},
  frame=single,
  framesep=3pt,
  xleftmargin=4pt,
  xrightmargin=4pt,
  keepspaces=true,
  columns=fullflexible,
  showstringspaces=false,
  breaklines=true,
  tabsize=2,
  escapeinside={(*@}{@*)},
  % accents UTF-8 dans le code (listings ne gere pas l'UTF-8 multi-octets)
  literate={é}{{\'e}}1 {è}{{\`e}}1 {ê}{{\^e}}1 {à}{{\`a}}1 {â}{{\^a}}1
           {ù}{{\`u}}1 {û}{{\^u}}1 {ô}{{\^o}}1 {î}{{\^i}}1 {ç}{{\c c}}1
           {É}{{\'E}}1 {À}{{\`A}}1 {€}{{\euro{}}}1 {…}{{\ldots}}1,
}

% ---------------------------------------------------------------
% CARD-STYLE BOXES
% ---------------------------------------------------------------
\tcbset{
  figcard/.style={colback=white, colframe=white, boxrule=0pt, arc=4pt,
    left=3pt, right=3pt, top=3pt, bottom=3pt, boxsep=0pt,
    halign=center, valign=center},
  datacard/.style={colback=qapanel, colframe=qaline, boxrule=0.4pt, arc=3pt,
    left=8pt, right=8pt, top=6pt, bottom=6pt, coltext=qatext,
    before upper={\color{qatext}}, halign=center},
  codebox/.style={colback=qapanelalt, colframe=qaline, boxrule=0.4pt, arc=3pt,
    left=8pt, right=8pt, top=5pt, bottom=5pt, listing only,
    listing options={basicstyle=\ttfamily\small\color{qaaccent},
      breaklines=true, showstringspaces=false, columns=fullflexible,
      keepspaces=true, frame=none, backgroundcolor={}}}
}
\newtcolorbox{figcard}[1][]{figcard,#1}
\newtcolorbox{datacard}{datacard}
\newtcblisting{codebox}{codebox}

% ---------------------------------------------------------------
% HANDY MACROS
% ---------------------------------------------------------------
% Mention d'elaboration (discrete : plus sombre que le texte secondaire)
\definecolor{qaaimention}{HTML}{4E5566}
\newcommand{\raimention}{{\tiny\color{qaaimention}Supports \'elabor\'es avec l'aide de Claude Code Opus 5.5}}
\newcommand{\hl}[1]{\alert{#1}}
\newcommand{\source}[1]{{\scriptsize\color{qamuted}#1}}

% Sommaire automatique en debut de section
\AtBeginSection[]{%
  \begin{frame}<beamer>{Plan du cours}
    \tableofcontents[currentsection]
  \end{frame}}

% Page de titre : \rtitleslide{numero du cours}{titre du cours}
% Les auteurs sont repris de \rauthors (modifiable avant \begin{document})
\newcommand{\rauthors}{Chouki Tibermacine}
\newcommand{\rtitleslide}[2]{%
\begin{frame}[plain,noframenumbering]
  \begin{center}
    \vspace{2mm}
    {\color{qaaccent}\rule{2.1cm}{1.2pt}}\\[6mm]
    {\Large\bfseries\color{qatext} R4.02}\\[1mm]
    {\Large\bfseries\color{qatext} Qualit\'e de d\'eveloppement}\\[5mm]
    {\large\color{qamuted} Cours #1}\\[1mm]
    {\large\color{qaaccent} #2}\\[6mm]
    \includegraphics[height=2.1cm]{iut_logo_transparent.png}

    \vspace{4mm}

    {\small\color{qatext} \rauthors}\\[1mm]
    {\scriptsize\color{qamuted} BUT 2 Informatique --- IUT de Vannes --- Universit\'e Bretagne Sud}\\[2mm]
    \raimention
  \end{center}
\end{frame}
}

% ---- Diapo de demonstration : \rdemo{numero}{titre}{contenu (items)}{fichier(s) a ouvrir} ----
\newcommand{\rdemo}[4]{%
\begin{frame}{D\'emo #1 --- #2}
  \begin{tikzpicture}[remember picture, overlay]
    \node[anchor=north east, fill=qaaccent2, text=qabg, font=\scriptsize\bfseries,
      rounded corners=3pt, inner sep=4pt] at ([xshift=-0.4cm, yshift=-0.35cm]current page.north east) {D\'EMO};
  \end{tikzpicture}
  \begin{columns}[T]
    \begin{column}{0.62\textwidth}
      \begin{itemize}
        #3
      \end{itemize}
    \end{column}
    \begin{column}{0.34\textwidth}
      \begin{datacard}
        \raggedright\scriptsize
        {\color{qamuted} Dans l'IDE}\\[1mm]
        #4
      \end{datacard}
    \end{column}
  \end{columns}
\end{frame}}

% ---- Diapo de fin « Questions ? » ----
\newcommand{\rcontact}{chouki.tibermacine@univ-ubs.fr}
\newcommand{\rquestionslide}{%
\begin{frame}{Questions ?}
  \begin{center}
    \vspace{2mm}
    \begin{tikzpicture}
      \node[regular polygon, regular polygon sides=6, draw=qaline, thick,
        minimum size=4.6cm, rotate=90] (hexin) {};
      \node[regular polygon, regular polygon sides=6, draw=qaaccent, dashed,
        thick, minimum size=6.2cm, rotate=90] (hexout) {};
      \node[font=\fontsize{92}{92}\selectfont\bfseries, color=qaaccent2] at (hexin) {?};
    \end{tikzpicture}\\[5mm]
    {\small\color{qamuted} \rcontact}
  \end{center}
\end{frame}}

\graphicspath{{../../latex/common/assets/}{img/}}

\title{User stories, tests d'acceptation et faisabilité}
\author[C. Tibermacine]{Chouki Tibermacine}
\date{}

% Gherkin (fichiers .feature), en francais
\lstdefinelanguage{Gherkin}{
  morekeywords={Fonctionnalité, Scénario, Plan, du, scénario, Exemples, Contexte,
    Étant, donné, Quand, Alors, Et, Mais, Feature, Scenario, Given, When, Then, And, But, Examples},
  sensitive=true,
  comment=[l]{\#},
  morestring=[b]"
}
\lstdefinestyle{gherkin}{language=Gherkin, basicstyle=\lstcode\scriptsize}
\lstdefinestyle{sortie}{language={}, basicstyle=\lstcode\scriptsize, keywordstyle={}, commentstyle={}, stringstyle={}}

\tikzset{
  etape/.style={rectangle, rounded corners=3pt, draw=qaaccent, thick, fill=qapanel, text=qatext,
    minimum height=0.7cm, minimum width=2.3cm, align=center, font=\scriptsize},
  arr/.style={-Latex, qamuted, thick}
}

\begin{document}

\rtitleslide{2}{User stories, tests d'acceptation et faisabilité}

\begin{frame}{Où en est-on ?}
  \begin{itemize}
    \item CM1 : \textbf{vérifier} que le code fait ce qu'il doit faire (tests unitaires de \texttt{Compte})
    \item Mais d'où viennent les règles R1 à R5 du TP1 ? Qui a décidé qu'un compte épargne est plafonné ?
    \item Aujourd'hui : \textbf{valider} qu'on construit le bon produit
      \begin{itemize}
        \item exprimer le besoin du client : les \textbf{user stories}
        \item le rendre vérifiable : les \textbf{critères d'acceptation}
        \item le rendre exécutable : les \textbf{tests d'acceptation} (Cucumber)
        \item avant de s'engager : la story est-elle \textbf{faisable} ?
      \end{itemize}
  \end{itemize}
  \vspace{3mm}
  \begin{datacard}
    \small \alert{besoin} $\rightarrow$ user story $\rightarrow$ critères d'acceptation $\rightarrow$ scénarios exécutables $\rightarrow$ code $\rightarrow$ recette
  \end{datacard}
\end{frame}

% ================================================================
\section{Méthodes agiles et user stories}

\begin{frame}{Les méthodes agiles en deux minutes}
  \begin{columns}[T]
    \begin{column}{0.5\textwidth}
      \begin{block}{Le Manifeste agile (2001)}
        \small
        \begin{itemize}
          \item Les individus et leurs interactions, plus que les processus et les outils
          \item Un logiciel qui fonctionne, plus qu'une documentation exhaustive
          \item La collaboration avec le client, plus que la négociation contractuelle
          \item L'adaptation au changement, plus que le suivi d'un plan
        \end{itemize}
      \end{block}
    \end{column}
    \begin{column}{0.46\textwidth}
      \begin{itemize}
        \item Développer par \textbf{petits incréments}, livrés et montrés au client régulièrement
        \item Méthodes les plus répandues : \textbf{Scrum} (organisation en \textit{sprints} de 1 à 4 semaines), \textbf{eXtreme Programming} (XP : pratiques techniques, dont les tests automatisés), Kanban
        \item Le besoin n'est pas figé dans un cahier des charges : il est découpé en petits morceaux, \alert{les user stories}, et précisé au fil de l'eau
      \end{itemize}
    \end{column}
  \end{columns}
\end{frame}

\begin{frame}{Les user stories dans Scrum}
  \begin{center}
  \begin{tikzpicture}[node distance=0.32cm, every node/.append style={minimum width=2.1cm}]
    \node[etape] (po) {\textit{Product owner}\\{\tiny\color{qamuted}porte le besoin}};
    \node[etape, right=of po, minimum width=2.5cm] (backlog) {\textit{Product backlog}\\{\tiny\color{qamuted}stories triées par priorité}};
    \node[etape, right=of backlog, minimum width=2.3cm] (sprint) {\textit{Sprint backlog}\\{\tiny\color{qamuted}stories choisies}};
    \node[etape, right=of sprint] (dev) {Sprint\\{\tiny\color{qamuted}1 à 4 semaines}};
    \node[etape, right=of dev, draw=qaaccent2] (incr) {Incrément\\{\tiny\color{qamuted}démontré, recetté}};
    \draw[arr] (po) -- (backlog);
    \draw[arr] (backlog) -- (sprint);
    \draw[arr] (sprint) -- (dev);
    \draw[arr] (dev) -- (incr);
    \draw[arr, qaaccent2] (incr.south) -- ++(0,-0.55) -| node[pos=0.25, below, font=\tiny, text=qamuted] {retours du client : nouvelles stories, priorités revues} (backlog.south);
  \end{tikzpicture}
  \end{center}
  \vspace{2mm}
  \begin{itemize}
    \item Le \textbf{product owner} représente le client : il écrit les stories et les ordonne
    \item L'\textbf{équipe de développement} les estime, les découpe et choisit celles qu'elle peut terminer pendant le sprint
    \item À la fin du sprint, les stories terminées sont \textbf{démontrées} et \textbf{validées} par rapport à leurs critères d'acceptation
  \end{itemize}
\end{frame}

\begin{frame}{Une user story}
  \begin{columns}[c]
    \begin{column}{0.5\textwidth}
      \begin{datacard}
        \raggedright
        \textbf{\color{qaaccent}En tant que} client de la banque,\\
        \textbf{\color{qaaccent}je veux} programmer un virement à une date future\\
        \textbf{\color{qaaccent}afin de} payer mon loyer sans y penser chaque mois.
      \end{datacard}
      \vspace{2mm}
      \begin{itemize}
        \item \textbf{Qui ?} un rôle (pas « l'utilisateur »)
        \item \textbf{Quoi ?} un besoin, pas une solution technique
        \item \textbf{Pourquoi ?} la valeur attendue : aide à arbitrer
      \end{itemize}
    \end{column}
    \begin{column}{0.46\textwidth}
      \begin{block}{Les 3C (Ron Jeffries, XP)}
        \begin{itemize}
          \item \textbf{Carte} : une phrase courte, qui tient sur un post-it
          \item \textbf{Conversation} : les détails se discutent entre le client et l'équipe
          \item \textbf{Confirmation} : les \alert{critères d'acceptation} disent quand la story est réalisée
        \end{itemize}
      \end{block}
      \source{Une story n'est pas une spécification complète : c'est une promesse de conversation.}
    \end{column}
  \end{columns}
\end{frame}

\begin{frame}{Une bonne story est INVEST}
  \begin{datacard}
    \small
    \begin{tabular}{@{}cl p{9.4cm}@{}}
      \textbf{\color{qaaccent}\Large I} & \textbf{Indépendante} & peut être développée et livrée sans attendre une autre story \\[0.5mm]
      \textbf{\color{qaaccent}\Large N} & \textbf{Négociable} & les détails restent ouverts à la discussion ; ce n'est pas un contrat \\[0.5mm]
      \textbf{\color{qaaccent}\Large V} & \textbf{Valeur} & apporte quelque chose à un utilisateur ou au client (pas « créer la table \texttt{Operation} ») \\[0.5mm]
      \textbf{\color{qaaccent}\Large E} & \textbf{Estimable} & l'équipe comprend assez bien la story pour estimer son coût \\[0.5mm]
      \textbf{\color{qaaccent}\Large S} & \textbf{Small} & petite : réalisable en quelques jours, bien moins qu'un sprint \\[0.5mm]
      \textbf{\color{qaaccent}\Large T} & \textbf{Testable} & on sait écrire des critères qui disent si elle est réalisée ou non \\
    \end{tabular}
  \end{datacard}
  \vspace{1mm}
  \source{Bill Wake, 2003. Une story trop grosse (\textit{epic}) se découpe en plusieurs stories.}
\end{frame}

\begin{frame}{Qu'est-ce qui ne va pas ?}
  \small
  \begin{enumerate}
    \item « En tant que développeur, je veux ajouter un index sur \texttt{Operation.numCompte}. »
      \uncover<2->{\\ \alert{$\rightarrow$ tâche technique, pas de valeur pour un utilisateur : c'est un moyen, pas un besoin}}
    \item « En tant que client, je veux gérer mes comptes. »
      \uncover<2->{\\ \alert{$\rightarrow$ trop grosse (\textit{epic}) : consulter, virer, programmer, clôturer\ldots{} à découper}}
    \item « En tant qu'utilisateur, je veux que l'application soit rapide. »
      \uncover<2->{\\ \alert{$\rightarrow$ non testable : rapide comment ? « la liste des comptes s'affiche en moins de 2 s »}}
    \item « En tant que client, je veux un bouton bleu en haut à droite qui ouvre une fenêtre modale. »
      \uncover<2->{\\ \alert{$\rightarrow$ décrit une solution, pas un besoin ; pas de « afin de » : pourquoi ce bouton ?}}
    \item « En tant que conseiller, je veux voir les comptes de mes clients, une fois que la story 12 sera finie. »
      \uncover<2->{\\ \alert{$\rightarrow$ dépendante d'une autre story : à réordonner ou à redécouper}}
  \end{enumerate}
\end{frame}

% ================================================================
\section{Critères d'acceptation}

\begin{frame}{Des règles\ldots{} et des exemples}
  \begin{columns}[T]
    \begin{column}{0.48\textwidth}
      \begin{block}{Règles (abstraites)}
        \small
        \begin{itemize}
          \item Un virement programmé est exécuté à la date choisie
          \item La date doit être dans le futur
          \item Si le solde est insuffisant le jour J, le virement est rejeté et le client est prévenu
        \end{itemize}
      \end{block}
    \end{column}
    \begin{column}{0.48\textwidth}
      \begin{exampleblock}{Exemples (concrets)}
        \small
        \begin{itemize}
          \item Solde 900 \euro{}, virement de 650 \euro{} programmé le 5 : le 5, le solde passe à 250 \euro{}
          \item Programmé « hier » : refusé à la saisie
          \item Solde 400 \euro{} le 5, virement de 650 \euro{} : rejeté, notification envoyée, solde inchangé
        \end{itemize}
      \end{exampleblock}
    \end{column}
  \end{columns}
  \vspace{3mm}
  \begin{itemize}
    \item Les exemples lèvent les ambiguïtés : \textit{« et si le 5 est un dimanche ? »} surgit dès qu'on écrit un exemple
    \item \textbf{Spécification par l'exemple} : client, développeurs et testeurs écrivent ensemble les exemples (les « trois amigos »)
    \item Chaque exemple deviendra un \textbf{test d'acceptation}
  \end{itemize}
\end{frame}

\begin{frame}[fragile]{Le format Gherkin : Étant donné / Quand / Alors}
\begin{lstlisting}[style=gherkin]
# language: fr
(*@\textcolor{qacodekw}{\textbf{Fonctionnalité:}}@*) Virement entre deux comptes
  En tant que client de la banque,
  je veux virer de l'argent entre mes comptes
  afin de gérer mon budget.
  (*@\textcolor{qacodekw}{\textbf{Scénario:}}@*) virement accepté
    (*@\textcolor{qacodekw}{\textbf{Étant donné}}@*) un compte "FR76-001" avec un solde de 500,00 €
    Et un compte "FR76-002" avec un solde de 100,00 €
    Quand le client vire 200,00 € de "FR76-001" vers "FR76-002"
    Alors le solde de "FR76-001" est de 300,00 €
    Et le solde de "FR76-002" est de 300,00 €

  (*@\textcolor{qacodekw}{\textbf{Scénario:}}@*) virement refusé, solde insuffisant
    (*@\textcolor{qacodekw}{\textbf{Étant donné}}@*) un compte "FR76-001" avec un solde de 50,00 €
    Quand le client vire 200,00 € de "FR76-001" vers "FR76-002"
    Alors le virement est refusé
    Et le solde de "FR76-001" est de 50,00 €
\end{lstlisting}
  \small
  \textbf{Étant donné} : le contexte (Arrange) \quad
  \textbf{Quand} : l'action (Act) \quad
  \textbf{Alors} : le résultat attendu (Assert) \\[1mm]
  Lisible par le client, mais \textbf{structuré} : chaque ligne pourra être exécutée.
\end{frame}

\begin{frame}[fragile]{Plusieurs exemples, un seul scénario}
\begin{lstlisting}[style=gherkin]
  (*@\textcolor{qacodekw}{\textbf{Plan du scénario:}}@*) virement avec confirmation au-delà de 1 000 €
    (*@\textcolor{qacodekw}{\textbf{Étant donné}}@*) un compte "FR76-001" avec un solde de 5000,00 €
    Quand le client vire <montant> € de "FR76-001" vers "FR76-002"
    Alors une confirmation est demandée : <confirmation>

    Exemples:
      | montant | confirmation |
      |  999,99 | non          |
      | 1000,00 | non          |
      | 1000,01 | oui          |
\end{lstlisting}
  \begin{itemize}
    \item L'équivalent d'un \textbf{test paramétré} JUnit, lisible par le client
    \item On y retrouve les \textbf{valeurs limites} du CM1 : la discussion « et à exactement 1\,000 \euro{} ? » a lieu \alert{avant} le code
    \item Bon critère d'acceptation : précis, vérifiable, indépendant de l'interface (pas « cliquer sur le bouton vert »)
  \end{itemize}
\end{frame}

\begin{frame}{Definition of Ready, Definition of Done}
  \begin{columns}[T]
    \begin{column}{0.48\textwidth}
      \begin{block}{Definition of Ready (DoR)}
        Quand une story peut-elle entrer dans un sprint ?
        \begin{itemize}
          \item elle respecte INVEST
          \item ses critères d'acceptation sont écrits et validés par le \textit{product owner}
          \item elle est estimée
          \item ses dépendances et risques sont connus
        \end{itemize}
      \end{block}
    \end{column}
    \begin{column}{0.48\textwidth}
      \begin{block}{Definition of Done (DoD)}
        Quand une story est-elle terminée ?
        \begin{itemize}
          \item tous ses critères d'acceptation passent
          \item code relu, tests unitaires et d'intégration écrits
          \item \textbf{aucune régression} : toute la suite de tests passe
          \item déployée en recette et validée par le client
        \end{itemize}
      \end{block}
    \end{column}
  \end{columns}
  \vspace{3mm}
  \begin{center}
    « Ça marche sur ma machine » n'est pas « terminé ». La DoD est commune à toute l'équipe.
  \end{center}
\end{frame}

\begin{frame}{Tests d'acceptation et recette}
  \begin{itemize}
    \item \textbf{Recette} : phase où le client vérifie que ce qui a été livré correspond à ce qui était demandé
      \begin{itemize}
        \item sur un environnement dédié, proche de la production (CM6)
        \item à l'aide des critères d'acceptation des stories
        \item conclue par un \textbf{procès-verbal de recette} : accepté, accepté avec réserves, refusé
      \end{itemize}
    \item Sans critères d'acceptation, la recette tourne à la négociation : « ce n'est pas ce que j'avais demandé »
    \item Si les critères sont \textbf{automatisés}, la recette porte sur ce qui ne s'automatise pas (ergonomie, cas imprévus) et on peut rejouer les scénarios à chaque livraison
  \end{itemize}
\end{frame}

% ================================================================
\section{Des scénarios exécutables : ATDD, BDD et Cucumber}

\begin{frame}{Écrire les tests d'acceptation avant le code}
  \begin{columns}[c]
    \begin{column}{0.52\textwidth}
      \begin{center}
      \begin{tikzpicture}
        \draw[qaaccent2, very thick, -Latex] (80:1.9) arc (80:-250:1.9);
        \draw[qaaccent, very thick, -Latex] (60:0.95) arc (60:-270:0.95);
        \node[font=\tiny, text=qaaccent2, align=center] at (0,1.55) {boucle d'acceptation};
        \node[font=\tiny, text=qaaccent2, align=center] at (-1.95,-1.95) {1. scénario\\rouge};
        \node[font=\tiny, text=qaaccent2, align=center] at (2.0,-1.95) {3. scénario\\vert};
        \node[font=\tiny, text=qaaccent, align=center] at (0,0) {2. tests unitaires\\rouge $\rightarrow$ vert\\(TDD)};
      \end{tikzpicture}
      \end{center}
    \end{column}
    \begin{column}{0.46\textwidth}
      \begin{itemize}
        \item \textbf{ATDD} (\textit{Acceptance Test-Driven Development}) :
          \begin{enumerate}
            \item écrire le scénario d'acceptation : il échoue
            \item développer, guidé par des tests unitaires
            \item le scénario passe : la story est réalisée
          \end{enumerate}
        \item \textbf{BDD} (\textit{Behaviour-Driven Development}, Dan North) : même idée, avec un vocabulaire métier partagé (Gherkin)
        \item Le scénario dit \alert{quand s'arrêter} : ni plus, ni moins que le besoin
      \end{itemize}
    \end{column}
  \end{columns}
\end{frame}

\begin{frame}{Cucumber : du texte au code}
  \begin{center}
  \begin{tikzpicture}[node distance=0.9cm]
    \node[etape, minimum width=3.2cm] (feat) {\texttt{virement.feature}\\{\tiny\color{qamuted}Gherkin, écrit avec le client}};
    \node[etape, right=of feat, minimum width=3.2cm] (steps) {\textit{Step definitions}\\{\tiny\color{qamuted}méthodes Java annotées}};
    \node[etape, right=of steps, minimum width=3.2cm] (app) {Code de l'application\\{\tiny\color{qamuted}\texttt{Compte}, \texttt{VirementService}\ldots}};
    \draw[arr] (feat) -- node[above, font=\tiny, text=qamuted] {chaque ligne} node[below, font=\tiny, text=qamuted] {cherche son étape} (steps);
    \draw[arr] (steps) -- node[above, font=\tiny, text=qamuted] {appelle} (app);
  \end{tikzpicture}
  \end{center}
  \vspace{2mm}
  \begin{itemize}
    \item \textbf{Cucumber} (Cucumber-JVM pour Java) lit les fichiers \texttt{.feature} et exécute, pour chaque ligne, la méthode Java dont l'expression correspond
    \item Les étapes sont \textbf{réutilisables} : « un compte \ldots{} avec un solde de \ldots{} » sert dans des dizaines de scénarios
    \item Exécuté par JUnit (\texttt{mvn test}), comme les autres tests : vert ou rouge, rapport HTML
    \item Les \texttt{.feature} deviennent une \textbf{documentation vivante} : elle ne peut pas être fausse, sinon les tests échouent
  \end{itemize}
\end{frame}

\begin{frame}[fragile]{Les \textit{step definitions}}
\begin{lstlisting}[basicstyle=\lstcode\scriptsize]
public class VirementSteps {
    private final ContexteBanque contexte;   // comptes en mémoire + VirementService

    public VirementSteps(ContexteBanque contexte) { this.contexte = contexte; }

    @Given("un compte {string} avec un solde de {bigdecimal} €")
    public void unCompte(String numero, BigDecimal solde) {
        contexte.ajouterCompte(new Compte(numero, solde));
    }
    @When("le client vire {bigdecimal} € de {string} vers {string}")
    public void vire(BigDecimal montant, String source, String cible) {
        contexte.virementService().virer(source, cible, montant);
    }
    @Then("le solde de {string} est de {bigdecimal} €")
    public void soldeEst(String numero, BigDecimal attendu) {
        assertThat(contexte.compte(numero).getSolde()).isEqualByComparingTo(attendu);
    }
}
\end{lstlisting}
  \source{\texttt{\{string\}}, \texttt{\{bigdecimal\}} : expressions Cucumber. Avec \texttt{\# language: fr}, les nombres s'écrivent à la française (\texttt{500,00}).}
\end{frame}

\rdemo{1}{Un scénario qui trouve un bug}{%
  \item Lire le fichier \texttt{.feature} comme le ferait le client
  \item Lancer : 4 scénarios verts, 1 exemple rouge : \texttt{1000,00 | non}
  \item Le client a dit « \textbf{au-delà} de 1\,000 \euro{} », le code dit \texttt{>=}
  \item Une étape inconnue : Cucumber propose le code à écrire
  \item Le rapport HTML, lisible par le client
}{\texttt{features/virement.feature}\\\texttt{acceptation/VirementSteps}\\\texttt{Banque}\\\texttt{target/cucumber.html}}

% ================================================================
\section{Critères de faisabilité}

\begin{frame}{Faisable ? Cinq questions (TELOS)}
  \small
  \begin{datacard}
    \begin{tabular}{@{}p{2.6cm}p{10.2cm}@{}}
      \toprule
      \textbf{\color{qaaccent}Technique} & Maîtrise-t-on les technologies ? Les données existent-elles ? Dépend-on d'un service externe (API) fiable, documenté, accessible ? \\
      \textbf{\color{qaaccent}Économique} & Le coût (développement, licences, hébergement, appels d'API payants) est-il justifié par la valeur apportée ? \\
      \textbf{\color{qaaccent}Légale} & Données personnelles (RGPD), réglementation du secteur (banque, santé), licences des bibliothèques utilisées \\
      \textbf{\color{qaaccent}Opérationnelle} & Les utilisateurs pourront-ils et voudront-ils s'en servir ? L'équipe a-t-elle les compétences ? Qui l'exploitera ? \\
      \textbf{\color{qaaccent}Temporelle} (\textit{Scheduling}) & Peut-on le livrer dans les délais ? Que faut-il abandonner sinon ? \\
      \bottomrule
    \end{tabular}
  \end{datacard}
  \vspace{2mm}
  \begin{itemize}
    \item Ces questions se posent pour un projet entier \textbf{et} pour chaque story, avant de s'engager
  \end{itemize}
\end{frame}

\begin{frame}{Estimer : points de story et \textit{planning poker}}
  \begin{columns}[T]
    \begin{column}{0.5\textwidth}
      \begin{itemize}
        \item On estime une \textbf{taille relative}, pas des heures : « deux fois plus gros que la story X »
        \item Échelle de Fibonacci : 1, 2, 3, 5, 8, 13, 20, 40, 100
          \begin{itemize}
            \item plus c'est gros, plus c'est incertain
            \item au-delà de 13 : story à découper
          \end{itemize}
        \item \textbf{Planning poker} : chacun choisit une carte en secret, on révèle ensemble, les extrêmes s'expliquent, on revote
      \end{itemize}
    \end{column}
    \begin{column}{0.46\textwidth}
      \begin{block}{Vélocité}
        \begin{itemize}
          \item Nombre de points terminés par sprint, mesuré sur les sprints passés
          \item Backlog de 120 points, vélocité de 20 : environ 6 sprints
          \item Permet de dire tôt : « tout ne tiendra pas dans les délais »
        \end{itemize}
      \end{block}
      \begin{alertblock}{Le « ? »}
        Personne ne sait estimer ? La story n'est pas faisable \textbf{en l'état} : il manque de l'information.
      \end{alertblock}
    \end{column}
  \end{columns}
\end{frame}

\begin{frame}{Les risques : probabilité $\times$ impact}
  \begin{columns}[c]
    \begin{column}{0.5\textwidth}
      \begin{center}
      \begin{tikzpicture}[x=1.25cm, y=0.95cm]
        \foreach \i in {0,1,2} \foreach \j in {0,1,2} {
          \pgfmathtruncatemacro{\s}{\i+\j}
          \ifnum\s<2 \def\c{qapanelalt}\else\ifnum\s<4 \def\c{qaaccent2!45!qapanel}\else\def\c{qaaccent!60!qapanel}\fi\fi
          \fill[\c, draw=qaline] (\i,\j) rectangle (\i+1,\j+1);
        }
        \node[font=\tiny, text=qamuted, rotate=90] at (-0.3,1.5) {probabilité};
        \node[font=\tiny, text=qamuted] at (1.5,-0.3) {impact};
        \node[font=\tiny, text=qatext, align=center] at (0.5,2.5) {A};
        \node[font=\tiny, text=qatext, align=center] at (1.5,1.5) {B};
        \node[font=\tiny, text=qatext, align=center] at (2.5,0.5) {C};
        \node[font=\tiny, text=qatext, align=center] at (2.5,1.5) {D};
      \end{tikzpicture}
      \end{center}
    \end{column}
    \begin{column}{0.48\textwidth}
      \small
      \begin{itemize}
        \item[A] L'équipe découvre Cucumber : retard de quelques jours
        \item[B] L'API de taux de change gratuite limite le nombre d'appels
        \item[C] Fuite de données bancaires
        \item[D] Le service SMS change ses conditions et devient payant
      \end{itemize}
      \vspace{2mm}
      Pour chaque risque : \textbf{éviter} (changer de solution), \textbf{réduire} (tester, prévoir un plan B), \textbf{transférer} (contrat, assurance), \textbf{accepter} (et le surveiller)
    \end{column}
  \end{columns}
\end{frame}

\begin{frame}{Lever une incertitude : le \textit{spike}}
  \begin{columns}[T]
    \begin{column}{0.5\textwidth}
      \begin{itemize}
        \item \textbf{Spike} (XP) : une courte exploration \textbf{limitée dans le temps} pour répondre à une question précise
        \item Exemple : \textit{« Quelle API de taux de change utiliser ? Limites d'appels ? Format ? Disponibilité ? »}
          \begin{itemize}
            \item 2 heures, un petit programme jetable
            \item résultat : une \textbf{décision}, et une story devenue estimable
          \end{itemize}
        \item \textbf{Preuve de concept} (POC) : même idée à l'échelle d'un projet (« la technologie X tient-elle la charge ? »)
      \end{itemize}
    \end{column}
    \begin{column}{0.46\textwidth}
      \begin{block}{La testabilité, critère de faisabilité}
        \begin{itemize}
          \item Une story dont on ne sait pas écrire les critères d'acceptation n'est \textbf{pas prête}
          \item Signes d'alerte : « rapide », « intuitif », « sécurisé », « comme l'ancien système »
          \item Des règles qui se \textbf{contredisent} apparaissent souvent en écrivant les exemples
          \item Une dépendance externe impossible à simuler rend les tests impossibles
        \end{itemize}
      \end{block}
    \end{column}
  \end{columns}
\end{frame}

\begin{frame}{La grille d'analyse de faisabilité (TP2)}
  \scriptsize
  \begin{datacard}
    \setlength{\tabcolsep}{3pt}
    \begin{tabular}{@{}p{2.0cm}p{2.0cm}p{0.8cm}p{1.5cm}p{0.9cm}p{1.7cm}p{1.1cm}p{1.6cm}@{}}
      \toprule
      \textbf{Story} & \textbf{Technique} & \textbf{Points} & \textbf{Dépendances} & \textbf{Légal} & \textbf{Testable ?} & \textbf{Risques} & \textbf{Décision} \\
      \midrule
      Virement programmé & maîtrisé (planificateur Spring) & 5 & horloge & -- & oui (exemples écrits) & faible & \textbf{\color{qaaccent}prête} \\[1mm]
      Conversion de devises & API externe inconnue & ? & API de taux & -- & partiellement & B & \textbf{\color{qaaccent2}spike} \\[1mm]
      « L'appli doit être sûre » & ? & ? & ? & RGPD & non & C & \textbf{\color{qaaccent2}reformuler} \\
      \bottomrule
    \end{tabular}
  \end{datacard}
  \vspace{2mm}
  \normalsize
  \begin{itemize}
    \item Décisions possibles : \textbf{prête}, \textbf{spike} d'abord, \textbf{reformuler} ou \textbf{découper}, \textbf{abandonner}
    \item Une décision se \textbf{justifie} : c'est ce qui sera évalué en TP, plus que la case cochée
  \end{itemize}
\end{frame}

\begin{frame}{À retenir}
  \footnotesize
  \begin{itemize}
    \item Les user stories viennent des \textbf{méthodes agiles} : le besoin est découpé en petits morceaux, ordonnés dans le \textit{product backlog} et précisés au fil des sprints
    \item « En tant que\ldots{} je veux\ldots{} afin de\ldots{} » ; une bonne story est \textbf{INVEST} ; carte, conversation, \textbf{confirmation}
    \item Les critères d'acceptation s'écrivent avec des \textbf{exemples} concrets (Étant donné / Quand / Alors), qui révèlent les ambiguïtés et les valeurs limites
    \item Avec Cucumber, ces exemples deviennent des \textbf{tests d'acceptation exécutables}, écrits \textbf{avant} le code (ATDD / BDD)
    \item Avant de s'engager : faisabilité \textbf{technique, économique, légale, opérationnelle, temporelle} ; estimer, évaluer les risques, lever les incertitudes par un \textit{spike}
    \item Une story non testable n'est pas prête
  \end{itemize}
  \vspace{0mm}
  \begin{exampleblock}{TP2}
    Rédiger des stories et leurs critères, analyser leur faisabilité, puis automatiser les scénarios avec Cucumber.
  \end{exampleblock}
\end{frame}

\rquestionslide

\end{document}
